\FloatBarrier
\subsection{Correctness and Scope}
\label{sec:correctness}
Synthesis proposes a construction for the replacement; verification must
establish that the construction preserves the original computation. A mapping
query can locate axes without establishing that attributes are evaluable,
operations legal, or computed tensors equal.

\FloatBarrier
\subsubsection{Checking the Completed Rewrite}
After synthesis produces $f$, the verifier substitutes its expressions into
the RHS and checks the completed rule in a fresh context. It reconstructs
tensor meanings and coordinate relationships rather than taking intermediate
synthesis conclusions as established facts. Table~\ref{tab:verification}
lists the checks.
\begin{table}[htbp]
\caption{Checks required for a completed semantic rewrite.}
\label{tab:verification}
\begin{tabularx}{\linewidth}{@{}lY@{}}\toprule
Check & Required fact\\\midrule
Attribute computation & Constructors depend only on $s$, are evaluable whenever the rule applies, and leave no attribute holes.\\
RHS legality & Constructed attributes satisfy operation definitions, including valid position lists, paired sizes, reduction domains, and access bounds.\\
Value equivalence & Normalized source and completed replacement agree in coefficients, indexed reads, and reduction domains.\\
Output agreement & Both sides have the same ordered output axes and shape, with compatible dtype as fixed by the contract.\\\bottomrule
\end{tabularx}
\end{table}
Read matching against a skeleton assists coordinate construction; it does not
prove value equivalence. That check uses the fully lowered RHS with actual
attributes and reduction axes. The verifier can assume source legality and
the declared guard. Any extra condition needed by construction must follow
from those assumptions and cannot be introduced merely to make checking
succeed. Explicit attribute programs undergo the same semantic checking.
The outcome is an accepted completed rule or an explanation of what could
not be established.

\FloatBarrier
\subsubsection{Correctness Guarantee}
\label{sec:guarantee}
The intended guarantee of successful semantic verification is the contract
in Equation~\eqref{eq:contract} for every legal guarded source instance,
rather than only particular shapes or test values. Specifically, for every
legal $s$ satisfying $\operatorname{Guard}(s)$, $f(s)$ should be evaluable,
make the RHS legal, and satisfy
\begin{equation}
 \forall X\text{ compatible with }s.\quad
 \operatorname{LHS}(s,X)=\operatorname{RHS}(s,f(s),X).
\end{equation}
Equality includes ordered output shape and element values, with element types
fixed as in Section~\ref{sec:problem}. It is relative to the selected arithmetic
policy; algebraic reassociation alone does not justify equality under strict
floating-point evaluation.

This is a correctness contract and intended guarantee, rather than a soundness
theorem established in this draft. Its assumptions include correct operation
definitions, sound primitive metadata and index checks, admissible coordinate
frames, and faithful evaluation of partial attribute programs. Checked
arithmetic, coordinate transport, and guard evaluation are implementation
obligations. A fresh context separates checking from construction, but both
paths share semantic infrastructure; this separation alone does not establish
implementation soundness. The current C++ implementation is not claimed to
have been formally proved sound. We present the verification explanation and
its assumptions without introducing a soundness proof.

\FloatBarrier
\subsubsection{Scope and Limitations}
Failure to establish a rule does not establish that the rule is incorrect.
The method is incomplete and accepts semantic rewrites only when supported
reasoning establishes all required conditions. Ambiguous read correspondence,
unsupported coupled or inverse attribute equations, executable index or guard
computations beyond the checker's reasoning fragment, unavailable arithmetic
laws, and bounded normalization or coordinate work can all prevent acceptance.
For example, several source-coordinate assignments for a replacement read
must remain unresolved unless the checker can establish agreement.

Diagnostics should distinguish malformed rules, unsupported or unresolved
reasoning, and exhausted proof-work limits. TEPL does not claim to find every
valid attribute program or the weakest sufficient guard. The intended tensor
equivalence guarantee applies to semantically verified rules. Structural rules
receive structural checks and do not carry that guarantee. Acceptance records
success of the supported checks under stated conditions; rejection as unproved
is not by itself a counterexample.
